\documentclass[11pt]{article}
% Shared preamble of the RMC kernel write-ups (% Shared preamble of the RMC kernel write-ups (% Shared preamble of the RMC kernel write-ups (\input{preamble} after \documentclass).
\usepackage[margin=1in]{geometry}
\usepackage{amsmath,amssymb,bm}
\usepackage{graphicx}
\usepackage{booktabs}
\usepackage{hyperref}

\newcommand{\dd}{\,\mathrm{d}}
\newcommand{\Ein}{E}
\newcommand{\Eout}{E'}
\newcommand{\Ecm}{E_{c}}
\newcommand{\mucm}{\mu_{c}}
\newcommand{\mulab}{\mu_0}

% Common closing section: how the verification figures are produced
\newcommand{\verificationmethod}{%
Each figure is produced by \texttt{verification/verify\_laws.py}. For an incident
energy $\Ein$ along $+z$, $10^6$ collisions are sampled with MC/DC's own reaction
sampler (\texttt{sample\_elastic\_scattering}, \texttt{sample\_inelastic\_scattering}
or \texttt{sample\_fission}), and every emitted neutron is histogrammed in
$48 \times 24$ bins of $(\Eout, \mulab)$, with $\mulab$ the lab cosine to $+z$. The
inverted kernel used by RMC is integrated over the same bins (Gauss--Legendre panels;
two-body lines are split where $\mulab^*(\Eout)$ crosses a cosine edge), multiplied by
the number of collisions, and compared with a $\chi^2$ test over bins with at least 5
expected counts (the rest pooled into one bin). The tables list $\chi^2/\mathrm{dof}$,
its $p$-value, the emitted neutrons per collision (sampled vs.\ the integral of the
inverted kernel), and the largest relative change of a bin integral when the
quadrature panels are doubled. Left panels: outgoing-energy marginal; middle:
lab-cosine marginal; right: per-bin $z = (\text{sampled} - \text{inverted}) /
\sqrt{\text{inverted}}$.}
 after \documentclass).
\usepackage[margin=1in]{geometry}
\usepackage{amsmath,amssymb,bm}
\usepackage{graphicx}
\usepackage{booktabs}
\usepackage{hyperref}

\newcommand{\dd}{\,\mathrm{d}}
\newcommand{\Ein}{E}
\newcommand{\Eout}{E'}
\newcommand{\Ecm}{E_{c}}
\newcommand{\mucm}{\mu_{c}}
\newcommand{\mulab}{\mu_0}

% Common closing section: how the verification figures are produced
\newcommand{\verificationmethod}{%
Each figure is produced by \texttt{verification/verify\_laws.py}. For an incident
energy $\Ein$ along $+z$, $10^6$ collisions are sampled with MC/DC's own reaction
sampler (\texttt{sample\_elastic\_scattering}, \texttt{sample\_inelastic\_scattering}
or \texttt{sample\_fission}), and every emitted neutron is histogrammed in
$48 \times 24$ bins of $(\Eout, \mulab)$, with $\mulab$ the lab cosine to $+z$. The
inverted kernel used by RMC is integrated over the same bins (Gauss--Legendre panels;
two-body lines are split where $\mulab^*(\Eout)$ crosses a cosine edge), multiplied by
the number of collisions, and compared with a $\chi^2$ test over bins with at least 5
expected counts (the rest pooled into one bin). The tables list $\chi^2/\mathrm{dof}$,
its $p$-value, the emitted neutrons per collision (sampled vs.\ the integral of the
inverted kernel), and the largest relative change of a bin integral when the
quadrature panels are doubled. Left panels: outgoing-energy marginal; middle:
lab-cosine marginal; right: per-bin $z = (\text{sampled} - \text{inverted}) /
\sqrt{\text{inverted}}$.}
 after \documentclass).
\usepackage[margin=1in]{geometry}
\usepackage{amsmath,amssymb,bm}
\usepackage{graphicx}
\usepackage{booktabs}
\usepackage{hyperref}

\newcommand{\dd}{\,\mathrm{d}}
\newcommand{\Ein}{E}
\newcommand{\Eout}{E'}
\newcommand{\Ecm}{E_{c}}
\newcommand{\mucm}{\mu_{c}}
\newcommand{\mulab}{\mu_0}

% Common closing section: how the verification figures are produced
\newcommand{\verificationmethod}{%
Each figure is produced by \texttt{verification/verify\_laws.py}. For an incident
energy $\Ein$ along $+z$, $10^6$ collisions are sampled with MC/DC's own reaction
sampler (\texttt{sample\_elastic\_scattering}, \texttt{sample\_inelastic\_scattering}
or \texttt{sample\_fission}), and every emitted neutron is histogrammed in
$48 \times 24$ bins of $(\Eout, \mulab)$, with $\mulab$ the lab cosine to $+z$. The
inverted kernel used by RMC is integrated over the same bins (Gauss--Legendre panels;
two-body lines are split where $\mulab^*(\Eout)$ crosses a cosine edge), multiplied by
the number of collisions, and compared with a $\chi^2$ test over bins with at least 5
expected counts (the rest pooled into one bin). The tables list $\chi^2/\mathrm{dof}$,
its $p$-value, the emitted neutrons per collision (sampled vs.\ the integral of the
inverted kernel), and the largest relative change of a bin integral when the
quadrature panels are doubled. Left panels: outgoing-energy marginal; middle:
lab-cosine marginal; right: per-bin $z = (\text{sampled} - \text{inverted}) /
\sqrt{\text{inverted}}$.}

\usepackage{amsthm}

\newtheorem{proposition}{Proposition}
\newcommand{\tpsi}{\tilde\psi}
\newcommand{\teps}{\tilde\epsilon}
\newcommand{\Sbar}{\bar S}
\newcommand{\Sigt}{\Sigma_t}

\title{A self-shielded (collision-density) energy trial space for RMC\\
\large design note, not yet implemented}
\date{}

\begin{document}
\maketitle

\noindent Builds on \texttt{linear\_energy\_basis.tex} (notation, the Legendre
polynomials $P_a(x_g(E))$ on bin $g$, and the slope tally) and \texttt{rmc\_math.tex}
(the residual and the two-phase schedule). $E'$ is the incident and $E$ the outgoing
energy.

\section{Motivation}

Two limitations of the current energy trial spaces can share one cause, the variation of
$\Sigt(E)$ \emph{inside} a bin.
\begin{enumerate}
  \item \emph{Stall and bias.} The collision residual $c(E) - \Sigt(E)\tpsi(E)$
  contains the in-bin shape of $\Sigt$, which no polynomial $\tpsi$ removes. Iterated
  with the pointwise $\Sigt$, RMC stalls at the noise of that term; iterated with the
  projected removal, a convergent iteration instead reaches a Galerkin fixed point
  that can exhibit the Jensen-type underestimate described in \texttt{rmc\_math.tex}
  under its constant-source assumption (O-16, 1--10\,MeV: 15--27\%
  in the 2.35 and 1.65\,MeV bins for the constant basis).
  \item \emph{Overshoot.} The linear basis removes most of that bias (1.9\% rms against
  4.6\% on O-16), but a straight line fitted to a flux with a narrow dip inside the bin
  pivots about the average and overshoots at the bin ends.
\end{enumerate}

\section{Exact factorization and the smoothness assumption}

In a nonfissioning infinite medium (0D) with unit-yield scattering, with
$F=\Sigt\phi$ the collision density and a normalized scalar scattering kernel $f$,
\begin{equation}
  F(E) = q(E) + \int \dd E'\, \frac{\Sigma_s(E')}{\Sigt(E')}\,F(E')\,f(E'\to E),
  \label{eq:collision}
\end{equation}
and $\phi = F/\Sigt$ exactly wherever $\Sigt>0$. Smoothness of $F$ is an additional
assumption. For classical elastic scattering on a nuclide of mass ratio $A$, the
energy support is $[\alpha E',E']$, $\alpha=((A-1)/(A+1))^2$; broad redistribution
can smooth a resonance when $\Gamma\ll(1-\alpha)E$, the source is smooth, and the
scattering ratios and angular law do not introduce comparable structure [1--3].
An integral over $E'$ alone does not guarantee this. For example, an
energy-preserving kernel gives $F=q\Sigt/\Sigma_a$ in a pure absorption/scattering
balance and can retain sharp cross-section structure.
The narrow-resonance approximation treats the numerator as slowly varying;
the algebraic factorization itself remains exact even when that approximation fails.

Possible nonsmooth features include the edges of a source, the
first-collision edges $\alpha E_0$ of a line source, and inelastic thresholds. These
can cause kinks or jumps in $\phi$ too. Aligning bin edges helps represent such
features, but a line source is distributional and needs explicit source treatment.

\section{Trial and test spaces}

Let $w(E) = 1/\Sigma_{t,m}(E)$ be the inverse total cross section of the material $m$
of the cell. Assume $0<\Sigma_{\min}\le\Sigma_{t,m}(E)\le\Sigma_{\max}<\infty$
on each bin, so the moments and trial functions are finite. This weight does not
apply in vacuum; very small cross sections can also impair conditioning.
On bin $g$, take
\begin{equation}
  \text{trial: } \tpsi(E) = w(E)\sum_{b=0}^{p}\tpsi_{gb}\,P_b(x_g(E)),\qquad
  \text{test: } P_a(x_g(E)),\quad a = 0..p,
\end{equation}
a Petrov--Galerkin pair. Then $\Sigt\tpsi = \sum_b \tpsi_{gb}P_b$ is a polynomial in
every bin: the trial space represents collision densities by polynomials.

\paragraph{Projection.} With the test functions, the projection $\Pi$ onto the trial
space matches the Legendre moments $t_a = \int_g\psi\,P_a\dd E$ of a function $\psi$:
\begin{equation}
  \sum_b B_{g,ab}\,(\Pi\psi)_{gb} = t_{ga},\qquad
  B_{g,ab} = \int_g w(E)\,P_aP_b\dd E .
  \label{eq:projection}
\end{equation}
$B_g$ is symmetric positive definite ($w>0$), $(p+1)\times(p+1)$ per bin. $\Pi$ is an
oblique projection: $\Pi v = v$ for every trial function $v$, and $\Pi\psi$ is the trial
function with the same moments as $\psi$. For $w$ constant on the bin it is the
Legendre ($L^2$) projection of \texttt{linear\_energy\_basis.tex}.

\paragraph{Tallies.} The moments $t_a$ are exactly what the existing track-length scores
estimate (\texttt{flux} for $a=0$, \texttt{flux-energy-slope} for $a=1$). The only new
step is the solve \eqref{eq:projection} for the correction $\teps$ after each transport
pass. (A Bubnov--Galerkin choice, test $= wP_a$, would need a new score
$w(E)P_a(x_g(E))$ per track and a different mass matrix; the Petrov--Galerkin choice
keeps the tallies and makes the removal trivial, below.)

\section{Galerkin equations and transfer moments}

Testing the transport equation with $P_a$ on bin $(g,j)$ (0D),
\begin{equation}
  \int_g P_a\,\Sigt\tpsi\dd E = \sum_b\tpsi_{gb}\int_g P_aP_b\dd E
  = \frac{\Delta E_g}{2a+1}\,\tpsi_{ga}
  \;\Longrightarrow\;
  \tpsi_{gja} = Q_{gja} + \Sbar_{gja},
  \label{eq:galerkin}
\end{equation}
i.e.\ in the normalization of \texttt{linear\_energy\_basis.tex} the removal matrix
$R_{g,ab} = \frac{2a+1}{\Delta E_g}\int_g \Sigt\,w\,P_aP_b\dd E = \delta_{ab}$ is the
identity, independent of the cross sections. The in-scatter projection uses moments
with the weight on the incident energy:
\begin{equation}
  M^{w,r}_{m,\,a'g'j',\,agj} = \int_{g'}\!\dd E'\,
  \frac{\sigma_r(E')}{\Sigma_{t,m}(E')}\,P_{a'}(x_{g'}(E'))
  \int_{j'}\!\dd\mu'\int_g\!\dd E\,P_a(x_g(E))\int_j\!\dd\mu\;
  \tau_r(E'\!\to E,\mu'\!\to\mu),
\end{equation}
$\Sbar_{gja} = \frac{2a+1}{\Delta E_g\Delta\mu_j}
\sum_{r\in m}\sum_{a'g'j'}N_rM^{w,r}_{m,\,a'g'j',\,agj}\tpsi_{g'j'a'}$,
with $N_r$ the nuclide's atom density. The material moment is the density-weighted
sum $M^w_m=\sum_{r\in m}N_rM^{w,r}_m$. The quadrature is
the existing one with one extra factor $w(E')$ in the incident integrand;
$\sigma_r/\Sigma_{t,m}\le 1/N_r$ is bounded and, inside a resonance of reaction $r$,
often smoother than $\sigma_r$ alone, but that is not guaranteed for arbitrary
mixtures and thresholds. The price is that $M^w$ depends on the material through
$w$: moments are computed and cached per material (composition and temperature) instead
of per nuclide.

\section{Residual}

\begin{proposition}[The collision residual is in the test space]
In the nonfissioning 0D formulation, with projected $Q$ and $\Sbar$, the weighted
trial space and pointwise $\Sigt(E)$ give
\[
  r_c(E) = Q(E) + \Sbar(E) - \Sigt(E)\tpsi(E)
  = \sum_a\big(Q_{gja} + \Sbar_{gja} - \tpsi_{gja}\big)P_a(x_g(E))
\]
on every bin: a polynomial of degree $p$, which vanishes identically at the Galerkin
fixed point \eqref{eq:galerkin}.
\end{proposition}

\noindent Consequences:
\begin{itemize}
  \item The pointwise and polynomial-projected removal coincide. No bin-averaged
  removal is needed, although omitting the scattering remainder still modifies the
  full operator. This eliminates that removal-shape mismatch; it does not eliminate
  all Galerkin error or imply convergence.
  \item If $q=Q$, angular remainders are absent, and there is no fission or streaming,
  the remaining residual at the Galerkin fixed point is
  $r_s=\mathcal S\tpsi-\Sbar$. For a sufficiently smooth scattering source with
  bounded derivatives, its polynomial approximation error is
  $O(\Delta E^{p+1})$. A lower residual and stall level are plausible benefits,
  not guarantees in every varying-cross-section bin. Tally variance also depends
  on sampling and transport response. General problems retain $q-Q$, fission,
  angular, spatial and boundary remainders as applicable.
  \item $|r_c|$ masses are exact and cheap ($|$polynomial$|$ on each bin, no
  cross-section grid).
\end{itemize}

\paragraph{Edge residual (1D).} Across a face between cells of the same material,
$D(E) = w(E)\sum_b(\tpsi_{R,b}-\tpsi_{L,b})P_b$; its mass $\int_g|D|\dd E$ needs the
cross-section grid (roots of a polynomial times a positive $w$ are the polynomial's
roots, so the integral splits at those and is a weighted polynomial integral). Between
different materials $w$ jumps, so $D$ does not vanish even for equal coefficients:
continuity in $z$ (planned) must then be imposed on $\tpsi$ itself, not on the
coefficients.

\paragraph{Scattering correction.} $T(E,\mu) = \int\dd E'\,\tpsi(E')\,U(E'\to E,\mu)$
now has the resonance structure of $w(E')$ in its integrand, so the incident-energy
integral of \texttt{in\_scatter\_density} must split at the cross-section grid points
in the kinematic window, as well as at trial and support boundaries, to resolve
that structure. This is a quadrature requirement for the proposed method, not
a claim that the current implementation already satisfies it. Pilot proposals
also require the signed-source support safeguards described in \texttt{rmc\_math.tex}.

\section{Expected accuracy and cost}

The following comparison concerns the nonfissioning 0D model with an exactly
represented source and no unresolved angular dependence; it is not a measured
performance comparison.

\begin{center}
\begin{tabular}{lll}
\toprule
 & Linear basis & Self-shielded linear basis \\
\midrule
Represented in each bin & $\phi$ by a line & $F = \Sigt\phi$ by a line \\
Collision residual at the fixed point & in-bin shape of $\Sigt\tpsi$ & 0 \\
Irreducible residual & $\Sigt$ shape $+$ $r_s$ & $r_s$ only \\
Transport per iteration & --- & unchanged (same scores) \\
Projection & diagonal & $2\times2$ solve per bin \\
Transfer moments & per nuclide & per material \\
Correction sampler & smooth $\tpsi(E')$ & split at the xs grid \\
\bottomrule
\end{tabular}
\end{center}

\noindent Where $\Sigt$ is constant on a bin, $w$ is constant and the self-shielded
basis spans the same functions as the linear basis. The A = 1 absorber (constant
cross sections) should therefore reproduce the same reconstructed fixed point.
Identical random seeds alone do not guarantee identical stochastic iterates if
proposal normalization, coefficient scaling or stopping criteria differ.

\section{Where the factorization weakens}

\begin{itemize}
  \item \emph{Leakage.} The scalar balance is
  $\Sigt\phi=q_0+\mathcal S_0\phi-\partial_zJ$, with
  $\phi=\int_{-1}^1\psi\dd\mu$, $J=\int_{-1}^1\mu\psi\dd\mu$, and $q_0$
  the angle-integrated source:
  in thin or boundary regions the flux does not dip as $1/\Sigt$ (the dip is partly
  filled by neutrons streaming in). An equivalence-theory-inspired option is the
  weight $w = 1/(\Sigt + \Sigma_e)$ with an escape cross section $\Sigma_e$ (Wigner's
  rational approximation for a lump) [2,3]. A geometry-dependent escape estimate
  is needed; $\Sigma_e=1/\bar\ell$ is an isolated-lump approximation, not a general
  leakage law. The weight may be chosen per region, but then
  $\Sigt\tpsi=[\Sigt/(\Sigt+\Sigma_e)]\sum_b\tpsi_bP_b$ is generally no longer
  polynomial. The exact-removal simplification does not carry over.
  \item \emph{Wide resonances.} Where $\Gamma \gtrsim (1-\alpha)E$ (heavy nuclides, low
  energies), the collision density can develop its own resonance structure [2];
  $F$ is then not smooth through the resonance and the polynomial part must represent
  that dip. The finite-dimensional Galerkin approximation is not exact; exact
  full-residual RMC can still recover the physical solution's projection in expectation.
  \item \emph{Upscatter and thermal energies.} Free-gas and thermal scattering broaden
  $f$, but thermal kernels and source structure require their own smoothness analysis.
  The algebraic factorization remains valid wherever $\Sigt>0$.
\end{itemize}

\section{Relation to existing methods}

\begin{itemize}
  \item \emph{Narrow-resonance weighting and Bondarenko self-shielding} [1--3] use
  shapes proportional to $1/(\sigma_t+\sigma_0)$ in microscopic units, or
  $1/(\Sigma_t+\Sigma_0)$ with consistent macroscopic units, as
  \emph{weighting functions} for multigroup constants. The multigroup unknowns
  are integrated group fluxes, rather than a resolved in-group shape. Here a
  related shape is the
  \emph{trial function} of a Petrov--Galerkin discretization, and the Monte Carlo
  correction (RMC), when its source and transport are unbiased, recovers the
  physical solution's projection in expectation.
  \item \emph{Subgroup and probability-table methods} [3] parametrize the flux by the
  value of $\Sigt$ inside a group; the weighted basis instead keeps a continuous energy
  variable with an analytic dependence on $\Sigt(E)$.
  \item \emph{Energy finite elements}: wavelet discontinuous Galerkin in energy with a
  $\Sigt$-orthogonalized basis for the resolved resonance range [4]; finite elements
  with discontiguous support (FEDS), which cluster representative spectra to define
  discontiguous energy elements and use spectral basis shapes [5]. Sorting only
  by $\Sigt$ with constant functions is a related band model, not the full published
  FEDS construction. The discrete generalized multigroup method expands the flux
  within coarse groups in discrete orthogonal polynomials [6].
  \item \emph{Monte Carlo functional expansion tallies} [7] estimate expansion
  coefficients such as $t_a$ with orthogonal polynomials; \eqref{eq:projection} adds a
  weighted mass matrix to such tallies.
  \item \emph{Residual / exponentially convergent Monte Carlo} has used piecewise
  constant and linear (continuous or discontinuous) trial spaces in space and angle
  [8--11], and piecewise-constant energy bins for continuous energy [12--14].
\end{itemize}
The search for this note did not identify an RMC implementation with this exact
$1/\Sigt$-weighted Petrov--Galerkin pair. This limited search does not establish
novelty or absence from the wider Galerkin transport literature. Its ingredients
(resonance weighting, energy finite elements and coefficient tallies) are established.

\section{Verification plan}

\begin{enumerate}
  \item A = 1 absorber (constant $\Sigt$): the reconstructed fixed point agrees
  with the linear basis up to numerical error; compare stochastic observables
  using their uncertainties rather than requiring bitwise equality for the same seeds.
  \item 0D O-16 and U-238 + H windows: the fixed point equals the deterministic solve of
  \eqref{eq:galerkin}; difference to SMC against the constant and linear bases; SMC
  tallied with the same moments $t_a$ (and finer bins) to compare like with like.
  \item Full-residual iterations: number of convergent iterations before the stall, and
  the stall level, against the linear basis.
\end{enumerate}

\begin{thebibliography}{99}
\bibitem{bondarenko} I.~I.~Bondarenko (ed.), \emph{Group Constants for Nuclear Reactor
  Calculations}, Consultants Bureau, New York (1964).
\bibitem{bell} G.~I.~Bell and S.~Glasstone, \emph{Nuclear Reactor Theory}, Van
  Nostrand Reinhold (1970) (resonance absorption, narrow- and wide-resonance
  approximations, equivalence relations).
\bibitem{stacey} W.~M.~Stacey, \emph{Nuclear Reactor Physics}, 2nd ed., Wiley-VCH
  (2007), Ch.~11, \emph{Resonance Absorption} (resonance approximations,
  equivalence relations and multiband self-shielding),
  \href{https://doi.org/10.1002/9783527611041.ch11}{doi:10.1002/9783527611041.ch11}.
\bibitem{letellier} R.~Le~Tellier, D.~Fournier and J.~M.~Ruggieri, ``A Wavelet-Based
  Finite Element Method for the Self-Shielding Issue in Neutron Transport,''
  \emph{Nucl.\ Sci.\ Eng.} \textbf{163}(1), 34--55 (2009),
  \href{https://doi.org/10.13182/NSE163-34}{doi:10.13182/NSE163-34}.
\bibitem{till} A.~T.~Till, \emph{Finite Elements with Discontiguous Support for Energy
  Discretization in Particle Transport}, Ph.D.\ thesis, Texas A\&M University (2015);
  A.~T.~Till, M.~L.~Adams and J.~E.~Morel, ``The Finite Element with Discontiguous
  Support Multigroup Method: Theory,'' M\&C 2015.
  Published treatment: A.~T.~Till, M.~L.~Adams and J.~E.~Morel,
  \emph{The Finite-Element with Discontiguous-Support Method},
  Nucl.\ Sci.\ Eng.\ \textbf{196}(1), 53--74 (2022),
  \href{https://doi.org/10.1080/00295639.2021.1932224}{doi:10.1080/00295639.2021.1932224}.
\bibitem{zhu} L.~Zhu and B.~Forget,
  \emph{A discrete generalized multigroup energy expansion theory},
  Nucl.\ Sci.\ Eng.\ \textbf{166}(3), 239--253 (2010),
  \href{https://doi.org/10.13182/NSE09-84}{doi:10.13182/NSE09-84};
  \emph{An energy recondensation method using the discrete generalized multigroup
  energy expansion theory}, Ann.\ Nucl.\ Energy \textbf{38}(8), 1718--1727 (2011),
  \href{https://doi.org/10.1016/j.anucene.2011.04.008}{doi:10.1016/j.anucene.2011.04.008}.
\bibitem{griesheimer} D.~P.~Griesheimer, W.~R.~Martin and J.~P.~Holloway,
  ``Convergence properties of Monte Carlo functional expansion tallies,''
  \emph{J.\ Comput.\ Phys.} \textbf{211}, 129 (2006).
\bibitem{peterson} J.~R.~Peterson, J.~E.~Morel and J.~C.~Ragusa,
  ``Exponentially-convergent Monte Carlo for the 1-D transport equation,'' M\&C 2013;
  J.~R.~Peterson, M.S.\ thesis, Texas A\&M University (May 2014),
  \href{https://oaktrust.library.tamu.edu/bitstreams/69ab2dae-3fef-4daf-96d0-2bd1896d81e3/download}
  {university copy}.
  Published formulation: \emph{Residual Monte Carlo for the One-Dimensional
  Particle Transport Equation}, SIAM J.\ Sci.\ Comput.\ \textbf{38}(6),
  B941--B961 (2016), \href{https://doi.org/10.1137/16M1057619}{doi:10.1137/16M1057619}.
\bibitem{franke} B.~C.~Franke, D.~E.~Bruss and J.~E.~Morel, ``Residual Monte Carlo with
  Discrete Scattering Angles in the 1-D Transport Equation,'' ANS Winter Meeting
  \emph{Trans.\ Am.\ Nucl.\ Soc.} \textbf{109}(1), 679--682 (2013),
  \href{https://www.ans.org/pubs/transactions/volume-109/}{ANS publication record}.
\bibitem{bolding} S.~R.~Bolding and J.~E.~Morel, ``Residual Monte Carlo Transport in
  Time with Consistent Low-Order Acceleration for 1D Thermal Radiative Transfer,''
  M\&C 2017,
  \href{https://www.kns.org/files/int_paper/paper/MC2017_2017_2/P102S02-03BoldingS.pdf}
  {conference paper}.
\bibitem{vermaak} J.~I.~C.~Vermaak and J.~E.~Morel, ``Transport Error Estimation Using
  Residual Monte Carlo,'' \emph{J.\ Comput.\ Phys.} (2022).
\bibitem{larsen-nse} M.~Larsen et al., ``A Residual Monte Carlo Algorithm for
  Continuous Energy Neutron Transport with Elastic Scattering,'' \emph{Nucl.\ Sci.\
  Eng.} \textbf{200}(3), 525--538 (2026),
  \href{https://doi.org/10.1080/00295639.2025.2495608}{doi:10.1080/00295639.2025.2495608}.
\bibitem{larsen-diss} M.~Larsen, Ph.D.\ dissertation, Oregon State University (2025).
\bibitem{larsen-1d} ``A Residual Monte Carlo Algorithm for 1D Continuous-Energy Neutron
  Transport,'' \emph{Nucl.\ Sci.\ Eng.} (2026), doi:10.1080/00295639.2026.2659992.
\end{thebibliography}

\end{document}
